\par\Needspace{8\baselineskip}\bigskip
\pdfbookmark[3]{Darstellung, 6}{pallas-II-article6}
\begin{center}6.\end{center}
\noindent\hspace*{\parindent}Wir bezeichnen den Bogen \(XN\), von \(X\) nach \(Y\) gemessen, mit \(\AscendingNode\) und die Summe dieses Bogens und des Bogens \(NL\), von \(N\) in Bewegungsrichtung gemessen, mit \(v\). Diese Summe nennen die Astronomen die Länge in der Bahn. Es sei \(\tilde\omega\) der Wert, den \(v\) annimmt, wenn \(L\) in \(P\) übergeht, also die Länge des Perihels. Dann ist \(v-\tilde\omega\) der Bogen \(PL\) oder die wahre Anomalie. Die Neigung der Bahn gegen den Ast des Großkreises \(XY\), der von \(N\) in rechtläufiger Richtung ausgeht, ist gleich \(ZM\). Wir bezeichnen sie mit \(i\). Damit lassen sich die vorhergehenden 15 Formeln \SourceLeaf{454}in der folgenden Form darstellen.
\[
\begin{array}{rl}
42)\quad&\begin{aligned}[t]x&=r(\cos\AscendingNode\mathbin{.}\cos(v-\AscendingNode)\\&\qquad-\sin\AscendingNode\mathbin{.}\sin(v-\AscendingNode)\mathbin{.}\cos i)\end{aligned}\\[5pt]
43)\quad&\begin{aligned}[t]y&=r(\sin\AscendingNode\mathbin{.}\cos(v-\AscendingNode)\\&\qquad+\cos\AscendingNode\mathbin{.}\sin(v-\AscendingNode)\mathbin{.}\cos i)\end{aligned}\\[5pt]
44)\quad&z=r\sin(v-\AscendingNode)\mathbin{.}\sin i
\end{array}
\]
\[
\begin{array}{rl}
45)\quad&A=\BareRoot{mp}\mathbin{.}\sin\AscendingNode\mathbin{.}\sin i\\[4pt]
46)\quad&B=-\BareRoot{mp}\mathbin{.}\cos\AscendingNode\mathbin{.}\sin i\\[4pt]
47)\quad&C=\BareRoot{mp}\mathbin{.}\cos i
\end{array}
\]
\[
\begin{array}{rl}
48)\quad&\begin{aligned}[t]Bz-Cy&=r\BareRoot{mp}\mathbin{.}(-\cos\AscendingNode\mathbin{.}\sin(v-\AscendingNode)\\&\qquad-\sin\AscendingNode\mathbin{.}\cos(v-\AscendingNode)\mathbin{.}\cos i)\end{aligned}\\[5pt]
49)\quad&\begin{aligned}[t]Cx-Az&=r\BareRoot{mp}\mathbin{.}(-\sin\AscendingNode\mathbin{.}\sin(v-\AscendingNode)\\&\qquad+\cos\AscendingNode\mathbin{.}\cos(v-\AscendingNode)\mathbin{.}\cos i)\end{aligned}\\[5pt]
50)\quad&Ay-Bx=r\BareRoot{mp}\mathbin{.}\cos(v-\AscendingNode)\mathbin{.}\sin i
\end{array}
\]
\[
\begin{array}{rl}
51)\quad&\begin{aligned}[t]CG-BH&=mpe(\cos\AscendingNode\mathbin{.}\cos(\tilde\omega-\AscendingNode)\\&\qquad-\sin\AscendingNode\mathbin{.}\sin(\tilde\omega-\AscendingNode)\mathbin{.}\cos i)\end{aligned}\\[5pt]
52)\quad&\begin{aligned}[t]AH-CF&=mpe(\sin\AscendingNode\mathbin{.}\cos(\tilde\omega-\AscendingNode)\\&\qquad+\cos\AscendingNode\mathbin{.}\sin(\tilde\omega-\AscendingNode)\mathbin{.}\cos i)\end{aligned}\\[5pt]
53)\quad&BF-AG=mpe\sin(\tilde\omega-\AscendingNode)\mathbin{.}\sin i
\end{array}
\]
\[
\begin{array}{rl}
54)\quad&\begin{aligned}[t]F&=e\BareRoot{mp}\mathbin{.}(-\cos\AscendingNode\mathbin{.}\sin(\tilde\omega-\AscendingNode)\\&\qquad-\sin\AscendingNode\mathbin{.}\cos(\tilde\omega-\AscendingNode)\mathbin{.}\cos i)\end{aligned}\\[5pt]
55)\quad&\begin{aligned}[t]G&=e\BareRoot{mp}\mathbin{.}(-\sin\AscendingNode\mathbin{.}\sin(\tilde\omega-\AscendingNode)\\&\qquad+\cos\AscendingNode\mathbin{.}\cos(\tilde\omega-\AscendingNode)\mathbin{.}\cos i)\end{aligned}\\[5pt]
56)\quad&H=e\BareRoot{mp}\mathbin{.}\cos(\tilde\omega-\AscendingNode)\mathbin{.}\sin i.
\end{array}
\]
\par Multiplizieren wir die Gleichungen 11, 12, 13 mit \(A,B,C\) und addieren sie, so finden wir
\[
x(BH-CG)+y(CF-AH)+z(AG-BF)=mp(r-p).
\]
\par Mithilfe der Gleichungen 27—29, 36—38 wird nun die linke Seite dieser Gleichung zu
\[
=-mper\cos PL=-mper\cos(v-\tilde\omega);
\]
also gilt
\[
\begin{array}{rl}57)\quad&r=\dfrac{p}{1+e\cos(v-\tilde\omega)}\mathpunct{\cdot}\end{array}
\]
\par Ebenso ergibt Gleichung 20 mithilfe von 27—29, 39—41
\[
er\BareRoot{mp}\mathbin{.}\cos\Pi L=\dfrac{pr\mathrm{d}r}{\mathrm{d}t}
\]
oder wegen
\[
\Pi L=v-\tilde\omega-90^{0}:
\]
\[
\begin{array}{rl}58)\quad&\dfrac{\mathrm{d}r}{\mathrm{d}t}=e\sin(v-\tilde\omega)\mathbin{.}\BareRoot{\dfrac{m}{p}}\mathpunct{\cdot}\end{array}
\]
\par Differenziert man Gleichung 57, so erhält man \(\dfrac{p\mathrm{d}r}{rr\mathrm{d}t}=e\sin(v-\tilde\omega)\mathbin{.}\dfrac{\mathrm{d}v}{\mathrm{d}t}\). \SourceLeaf{455}Verbindet man dieses Ergebnis mit Gleichung 58, so erhält man
\[
\begin{array}{rl}59)\quad&\dfrac{rr\mathrm{d}v}{\mathrm{d}t}=\BareRoot{mp}.\end{array}
\]
\par Man kann bemerken, dass \(\tfrac12rr\mathrm{d}v\) das Flächenelement ist, das während des Zeitelements um die Sonne überstrichen wird. Die Projektionen dieser Fläche auf die drei Grundebenen erhält man, indem man sie mit \(\cos XM\), \(\cos YM\), \(\cos ZM\) multipliziert. Also müssen \(\tfrac12A\mathrm{d}t\), \(\tfrac12B\mathrm{d}t\), \(\tfrac12C\mathrm{d}t\) diesen Projektionen gleich sein, was sich übrigens unmittelbar überprüfen lässt.
\par Schließlich bemerken wir, dass
\[
\begin{aligned}
r\sin(v-\tilde\omega)&=\BareRoot{\dfrac{p}{m}}\mathbin{.}\dfrac{r\mathrm{d}r}{e\mathrm{d}t}\\
&=\BareRoot{ap}\mathbin{.}\sin E=a\BareRoot{(1-ee)}\mathbin{.}\sin E
\end{aligned}
\]
\[
\begin{aligned}
r+r\cos(v-\tilde\omega)&=\dfrac{p-r(1-e)}{e}\\
&=\dfrac{p-a(1-e)(1-e\cos E)}{e}=a(1-e)(1+\cos E)
\end{aligned}
\]
\[
\begin{aligned}
r-r\cos(v-\tilde\omega)&=\dfrac{-p+r(1+e)}{e}\\
&=\dfrac{-p+a(1+e)(1-e\cos E)}{e}=a(1+e)(1-\cos E).
\end{aligned}
\]
\par Teilt man die erste dieser Gleichungen durch die zweite oder die dritte durch die erste, so erhält man
\[
\begin{array}{rl}60)\quad&\operatorname{tang}\tfrac12(v-\tilde\omega)=\BareRoot{\dfrac{1+e}{1-e}}\mathbin{.}\operatorname{tang}\tfrac12E.\end{array}
\]
\par\medskip\begin{center}\rule{.18\textwidth}{.3pt}\end{center}
